\thispagestyle{empty}
    \begin{minipage}{0.22\linewidth}
    \begin{tcolorbox}[colback=white, title=, width=1\linewidth]	
        BTS  \\
        Spécialité CIEL 
    \end{tcolorbox}
\end{minipage}\hfill
\begin{minipage}{0.55\linewidth}
    \begin{tcolorbox}[colback=white, title=, width=1\linewidth]	
        \begin{center}
            Lycée Victor Hugo~--~COLOMIERS \\
            Epreuve ponctuelle Mathématiques
        \end{center}
        
    \end{tcolorbox}
\end{minipage}\hfill
\begin{minipage}{0.23\linewidth}
    \begin{tcolorbox}[colback=white, title=, width=1\linewidth]	
        \begin{center}
            Lundi 18 Mai 2026 \\
    Durée 60 min
        \end{center}
        
    \end{tcolorbox}
\end{minipage}\hfill

\vspace{2cm}

\NomPrenom{}

\vspace{2cm}

\begin{center}
    
\emph{Un ordinateur doté d'un tableur et du logiciel Géogébra est mis à la disposition du candidat qui pourra également se servir de sa calculette scientifique.}

\bigskip

\emph{Aucun document n'est autorisé.}

\bigskip

\emph{L'usage du téléphone est interdit.}

\bigskip

\emph{La clarté du raisonnement et la qualité de la rédaction entreront pour une part importante dans l'appréciation des copies.}

\end{center}

\bigskip

\vspace{2cm}

Ce sujet comporte~\pageref{LastPage} pages numérotées de 1 à~\pageref{LastPage} et est constitué de 3 parties indépendantes.

\vspace{2cm}


S'il le juge nécessaire, il pourra préparer un support (feuille, tableur, document quelconque) pour présenter ses réponses lors de l'exposé oral. Il devra sauvegarder son travail sur la clé USB fournie. Dans tous les cas, cette clé USB devra être remise au jury à la fin de l'entretien.
\vspace{1cm}


Cette épreuve est une épreuve orale d'une durée de 1h35 maximum:
\begin{itemize}
    \item \textbf{Préparation 1 heure}: le candidat  préparera les réponses aux différentes questions. Il pourra répondre directement sur le sujet ou sur une feuille séparée.
    \item \textbf{Exposé 15 minutes maximum}: le candidat présentera les réponses aux différentes questions à l'aide d'un support de son choix (feuille, tableau, ordinateur, etc.). Le candidat pourra s'aider de son ordinateur pour présenter ses réponses.
    \item \textbf{Entretien 20 minutes maximum}: le jury posera des questions au candidat sur les différentes parties du sujet. Le candidat pourra s'aider de son ordinateur pour répondre aux questions du jury.
\end{itemize}


\vspace{2cm}

\emph{Le jury se réserve le droit de ne pas interroger le candidat sur toutes les questions du sujet et pourra choisir d'ajouter des questions supplémentaires en fonction des réponses du candidat lors de l'entretien.
}

\vfill

\newpage